\chapter{Conclusion}

The demand for transparency in machine learning environments can no longer be ignored. The ZKLR project successfully demonstrated that it is possible to bridge the gap between proprietary intelligence models and mathematical accountability.

In Phase 1, we defined a mathematical bridge converting abstract predictive equations into strict zero-knowledge arithmetic constraints, delivering a baseline prototype capable of executing private evaluations resulting in publicly verifiable, 584-byte constant proofs.

This report, encapsulating Phase 2, attacked the ultimate barrier: real-world computational deployability. By aggressively refactoring the constraints via truncated table approximations, generalizing dynamic features, and deploying asynchronous batch parallelism, latency was virtually eliminated. Throughput surged by a factor of $12.5\times$, achieving evaluation velocities of $0.160$ seconds per inference under stable, industry-aligned BN254 constraints.

As algorithms continue expanding their governance over human decisions, frameworks like ZKLR prove that cryptography can offer an elegant balance—simultaneously shielding valuable intellect frameworks alongside mathematically guaranteeing a user their right to truthful outputs.
